%! Principles of Experimental Design — Unit 1, Lesson 1.6 — Group Activity (Answer Key)
\documentclass[10pt]{article}
\usepackage{saar-article}
\usepackage{saar-key}   % pulls in -boxes; do NOT also load -boxes
\newcommand{\TallMath}[1]{$\displaystyle #1\rule[-1.4em]{0pt}{3.2em}$}

\begin{document}

\pageheader{Unit 1, Lesson 1.6}{Group Activity --- Answer Key}
% NAMESTRIP: no \namedateperiod here — mirrors the blank. Teacher-only prose goes
% in the lesson plan as \begin{teachernote}[Group Activity], never in a key.

\vspace{0.15in}

\begin{headlinebox}{frost}
\small
\textbf{Settle It --- Harbor Point Community Pool.} Last lesson the pool's
survey found that $70\%$ of morning swimmers sleep at least $7$ hours against
$45\%$ of evening swimmers --- but the members had chosen their own swim times,
and $75\%$ of the morning swimmers turned out to be retired. This month the
director settles the question. \textbf{$90$ members volunteer.} A random number
generator assigns \textbf{$45$ to a reserved $6$~a.m.\ slot} and \textbf{$45$ to
a reserved $7$~p.m.\ slot} for four weeks; everything else about the pool stays
the same. Every volunteer keeps a sleep log. At the end, \textbf{$27$} of the
morning group averaged at least $7$ hours of sleep, and \textbf{$18$} of the
evening group did.
\end{headlinebox}

\vspace{0.03in}

\boxguard[30]
\begin{tcolorbox}[colback=white, colframe=black!40, breakable,
    title=\textbf{Tier R --- Remediate},
    fonttitle=\bfseries, arc=2mm, left=3mm, right=3mm, top=1mm, bottom=1mm]
\small
\begin{enumerate}[leftmargin=1.6em, itemsep=2pt, topsep=3pt]

  \item Name the numbers in this experiment.

        Volunteers: \ans{$90$} \hfill Morning group: \ans{$45$} \hfill
        Evening group: \ans{$45$} \hfill Pool members in all: \ans{$1{,}500$}

  \item Nobody asked the volunteers what time they would prefer.

        \textbf{(a)} Is this an experiment? \ans{yes}

        \textbf{(b)} Who decided which group each volunteer ended up in?
        \ans{chance --- a random generator}

  \item Of the $45$ volunteers assigned to the morning slot, $27$ averaged at
        least $7$ hours of sleep. What percent is that?

        \begin{work}
          \dfrac{27}{45} &= 0.60 \\
                         &= 60\%
        \end{work}

  \item Of the $45$ assigned to the evening slot, $18$ did. What percent is
        that?

        \begin{work}
          \dfrac{18}{45} &= 0.40 \\
                         &= 40\%
        \end{work}

  \item Which group slept more, and by how many percentage points?
        \ansline{The morning group, by $60 - 40 = 20$ percentage points.}
\end{enumerate}
\end{tcolorbox}

\vspace{0.03in}

\boxguard
\begin{tcolorbox}[colback=white, colframe=black!40, breakable,
    title=\textbf{Tier A --- Approaching Proficiency},
    fonttitle=\bfseries, arc=2mm, left=3mm, right=3mm, top=1.5mm, bottom=1.5mm]
\small
\begin{enumerate}[leftmargin=1.6em, itemsep=3pt, topsep=3pt]

  \item Put the volunteers back together.

        \textbf{(a)} How many of the $90$ volunteers slept at least $7$ hours,
        and what percent is that?

        \begin{work}
          27 + 18 &= 45 \text{ volunteers} \\
          \dfrac{45}{90} &= 0.50 = 50\%
        \end{work}

        \textbf{(b)} About how many of all $1{,}500$ pool members does that
        predict?

        \begin{work}
          0.50 \times 1500 &= 750 \text{ members}
        \end{work}

  \item Name the parts of the design.

        Experimental units: \ans{the $90$ volunteers} \hfill Treatment:
        \ans{the $6$~a.m.\ slot}

        Control group: \ans{the $7$~p.m.\ group}

  \item In a few words each, say how this design meets the four principles.

        Comparison: \ans{two groups, morning against evening} \hfill
        Randomization: \ans{a generator made the split}

        Replication: \ans{$45$ volunteers in each group, not one} \hfill
        Control: \ans{same pool, same four weeks, same everything else}

  \item Watching gave a $25$-point gap. This experiment gives a $20$-point gap.
        Explain why the \emph{smaller} number is the one to trust.
        \ansline{Watching mixed the swim time with retirement, which inflated the gap. Random}
        \ansline{assignment spread the retired members evenly, so $20$ points is the swim time alone.}
\end{enumerate}
\end{tcolorbox}

\vspace{0.03in}

\boxguard
\begin{tcolorbox}[colback=white, colframe=black!40, breakable,
    title=\textbf{Tier E --- Extension},
    fonttitle=\bfseries, arc=2mm, left=3mm, right=3mm, top=1.5mm, bottom=1.5mm]
\small
\begin{enumerate}[leftmargin=1.6em, itemsep=4pt, topsep=3pt]

  \item Retirement was the lurking variable last lesson, so the director blocks
        on it: $30$ of the $90$ volunteers are retired and $60$ still work full
        time. She randomizes \emph{inside} each block.

        \begin{center}
        \renewcommand{\arraystretch}{1.35}
        \begin{tabularx}{0.93\linewidth}{@{} X c c c @{}}
        \toprule
        \textbf{Block} & \textbf{In block} & \textbf{To $6$~a.m.}
          & \textbf{To $7$~p.m.} \\
        \midrule
        Retired & $30$ & \ans{$15$} & \ans{$15$} \\
        Still working full time & $60$ & \ans{$30$} & \ans{$30$} \\
        \midrule
        Total & $90$ & \ans{$45$} & \ans{$45$} \\
        \bottomrule
        \end{tabularx}
        \end{center}

        Why block on retirement instead of trusting the random assignment to
        even it out? \ansline{Chance usually evens it out, but it can miss; blocking makes it certain.}

  \item Each change below breaks one principle. Name it, and say what could go
        wrong.

        \textbf{(a)} The director lets each volunteer pick the slot they prefer.
        \ansline{Randomization. The retired members pick mornings again, and we are back to 1.5.}

        \textbf{(b)} The morning group swims in the newly renovated lap pool;
        the evening group swims in the old one.
        \ansline{Control. The new pool is a confounding variable --- the gap may be the pool.}

  \item Could this experiment be run \emph{double-blind}? Say who could be kept
        from knowing which group a volunteer was in, and who could not.
        \ansline{No. A volunteer obviously knows whether they swam at $6$ a.m.\ or $7$ p.m. But the}
        \ansline{staff who read and score the sleep logs can be kept from knowing --- that is single-blind.}
\end{enumerate}
\end{tcolorbox}

\end{document}
